\documentclass[UTF8,a4paper,12pt]{ctexart}

\usepackage[a4paper,top=2.35cm,bottom=2.35cm,left=2.55cm,right=2.35cm,headheight=15pt]{geometry}
\usepackage{amsmath,amssymb,bm}
\usepackage{booktabs,longtable,tabularx,array,multirow}
\usepackage{graphicx,float,caption,subcaption}
\usepackage{xcolor}
\usepackage{listings}
\usepackage{tikz}
\usetikzlibrary{arrows.meta,positioning,shapes.geometric,fit,calc,backgrounds}
\usepackage{pgfplots}
\pgfplotsset{compat=1.18}
\usepackage{fancyhdr}
\usepackage{hyperref}
\usepackage{enumitem}
\usepackage{microtype}
\usepackage{lastpage}

\definecolor{deepblue}{RGB}{28,74,121}
\definecolor{lightblue}{RGB}{232,242,252}
\definecolor{deepgreen}{RGB}{28,112,83}
\definecolor{lightgreen}{RGB}{232,247,240}
\definecolor{deeporange}{RGB}{185,92,23}
\definecolor{lightorange}{RGB}{252,241,229}
\definecolor{codered}{RGB}{165,42,42}
\definecolor{codegray}{RGB}{95,99,104}
\definecolor{codebg}{RGB}{247,248,250}

\hypersetup{
  colorlinks=true,
  linkcolor=deepblue,
  citecolor=deepgreen,
  urlcolor=deepblue,
  pdftitle={作业4：SM3软件实现与优化},
  pdfauthor={学生}
}

\pagestyle{fancy}
\fancyhf{}
\fancyhead[L]{作业4：SM3 软件实现与优化}
\fancyhead[R]{ARM64 NEON / x86 AVX2 / AVX-512}
\fancyfoot[C]{第 \thepage\ 页，共 \pageref{LastPage} 页}
\renewcommand{\headrulewidth}{0.4pt}
\renewcommand{\footrulewidth}{0pt}

\setlength{\parindent}{2em}
\setlength{\parskip}{0.25em}
\setlist[itemize]{leftmargin=2.1em,itemsep=0.15em,topsep=0.2em}
\setlist[enumerate]{leftmargin=2.5em,itemsep=0.15em,topsep=0.2em}
\renewcommand{\arraystretch}{1.22}

\lstdefinestyle{cstyle}{
  language=C,
  basicstyle=\ttfamily\small,
  keywordstyle=\color{deepblue}\bfseries,
  commentstyle=\color{deepgreen},
  stringstyle=\color{codered},
  numberstyle=\tiny\color{codegray},
  numbers=left,
  stepnumber=1,
  numbersep=8pt,
  backgroundcolor=\color{codebg},
  frame=single,
  rulecolor=\color{black!18},
  breaklines=true,
  showstringspaces=false,
  tabsize=4,
  columns=fullflexible,
  keepspaces=true
}
\lstdefinestyle{shellstyle}{
  language=bash,
  basicstyle=\ttfamily\small,
  keywordstyle=\color{deepblue},
  commentstyle=\color{deepgreen},
  backgroundcolor=\color{codebg},
  frame=single,
  rulecolor=\color{black!18},
  breaklines=true,
  showstringspaces=false
}

\newcommand{\code}[1]{\texttt{#1}}
\newcommand{\backend}[1]{\texttt{#1}}
\newcommand{\vect}[1]{\bm{#1}}
\newcommand{\rotl}{\operatorname{ROTL}_{32}}

\newenvironment{keybox}[1]{%
  \begin{center}\begin{minipage}{0.94\textwidth}
  \colorbox{lightblue}{\parbox{0.965\textwidth}{\textbf{#1}}}\vspace{0.35em}\par
}{%
  \end{minipage}\end{center}
}

\begin{document}

\begin{titlepage}
  \centering
  \vspace*{1.0cm}
  {\zihao{-2}\bfseries 网络安全创新创业实践课实验报告\par}
  \vspace{1.4cm}
  {\zihao{1}\bfseries 作业4：SM3 软件实现与优化\par}
  \vspace{0.55cm}
  {\zihao{3}\color{deepblue}\bfseries 基于 SIMD 寄存器与通用寄存器混合的\par}
  \vspace{0.18cm}
  {\zihao{3}\color{deepblue}\bfseries ARM64 NEON、x86 AVX2 与 AVX-512 实现\par}
  \vspace{1.4cm}

  \begin{tikzpicture}[node distance=0.55cm]
    \node[draw=deepblue,rounded corners=4pt,fill=lightblue,minimum width=3.5cm,minimum height=1.2cm,align=center] (scalar) {标量参考实现\\正确性基线};
    \node[draw=deepgreen,rounded corners=4pt,fill=lightgreen,minimum width=3.5cm,minimum height=1.2cm,right=0.65cm of scalar,align=center] (avx2) {x86 AVX2\\8 路并行};
    \node[draw=deeporange,rounded corners=4pt,fill=lightorange,minimum width=3.5cm,minimum height=1.2cm,right=0.65cm of avx2,align=center] (avx512) {x86 AVX-512\\16 路并行};
    \node[draw=purple!70!black,rounded corners=4pt,fill=purple!7,minimum width=3.5cm,minimum height=1.2cm,below=0.75cm of avx2,align=center] (neon) {ARM64 NEON\\4 路并行};
    \draw[-{Latex[length=2mm]},thick] (scalar) -- (avx2);
    \draw[-{Latex[length=2mm]},thick] (avx2) -- (avx512);
    \draw[-{Latex[length=2mm]},thick] (scalar.south) |- (neon.west);
  \end{tikzpicture}

  \vfill
  \renewcommand{\arraystretch}{1.7}
  \begin{tabular}{>{\bfseries}p{3.2cm}p{8.2cm}}
    课程名称： & \hrulefill \\
    姓名： & \hrulefill \\
    学号： & \hrulefill \\
    班级： & \hrulefill \\
    指导教师： & \hrulefill \\
    完成日期： & 2026 年 7 月 27 日 \\
  \end{tabular}
  \vfill

  {\small\color{codegray}
  说明：报告中的“本人 x86 实测”采用已保存的 AMD Ryzen 5 5600U / WSL2 运行记录；
  AVX-512 数据来自独立 x86-64 验证环境。ARM64 代码已完成交叉编译与指令反汇编验证，
  但未在当前 x86 主机上伪造 ARM64 运行性能。}
\end{titlepage}

\pagenumbering{Roman}
\section*{摘要}
\addcontentsline{toc}{section}{摘要}
本实验围绕国家商用密码杂凑算法 SM3，完成了从标量基线到多架构 SIMD 批量优化的完整软件工程。标量实现严格对应消息填充、消息扩展和 64 轮压缩函数，用作统一正确性基线。在 x86-64 平台上，分别实现 AVX2 八路和 AVX-512 十六路多缓冲区（multi-buffer）后端；在 ARM64 平台上实现 NEON 四路后端。三种向量后端均采用结构数组（SoA）布局：一个向量寄存器保存多条独立消息的同一状态字或同一扩展字，使消息扩展、布尔函数、模加、置换函数和循环移位能够按 lane 并行执行。

工程重点实现了题目要求的“SIMD 寄存器与通用寄存器混合”。SIMD 寄存器负责 $A\sim H$、$W_j$、$W'_j$、$P_0/P_1$、$FF/GG$ 和轮内模加；通用寄存器负责消息指针、长度、分组索引、活跃 lane 位图、填充尾块及运行时后端调度。针对同一批次消息长度不一致的问题，AVX2 使用向量 blend，AVX-512 使用 opmask，NEON 使用 bit-select，只提交仍处于活跃状态的 lane。自动分发按“AVX-512 十六路 $\rightarrow$ AVX2 八路 $\rightarrow$ 标量”或“NEON 四路 $\rightarrow$ 标量”处理任意数量消息。

正确性方面，工程通过空消息、\code{abc}、\code{abcd} 重复 16 次等标准向量，覆盖 55/56、63/64、119/120 字节等填充边界，并完成随机差分测试、128 个文件的 OpenSSL 差分、ASan/UBSan 检查和目标指令反汇编验证。在本人 AMD Ryzen 5 5600U / WSL2 的已保存测试中，AVX2 在 64 B 至 1 MiB 消息上达到约 $6.28\sim 8.56$ 倍批量吞吐率加速；独立 AVX-512 验证环境在 4096 B 消息上达到 1843.98 MiB/s，相对同环境标量实现约 12.89 倍。实验说明，宽 SIMD 对多条独立消息的吞吐率提升显著，但该实现的目标不是降低单条短消息延迟。

\noindent\textbf{关键词：} SM3；SIMD；通用寄存器；AVX2；AVX-512；ARM64；NEON；多缓冲区；运行时分发

\clearpage
\tableofcontents
\clearpage
\pagenumbering{arabic}

\section{实验任务与目标}

\subsection{作业要求}
本作业要求完成 SM3 软件实现与优化，并满足以下核心约束：
\begin{enumerate}
  \item 实现正确的 SM3 消息填充、消息扩展和迭代压缩；
  \item 采用 SIMD 寄存器和通用寄存器混合的优化方式；
  \item 覆盖 ARM64 指令集和 x86 指令集；
  \item x86 端至少实现 AVX2 与 AVX-512 两类宽向量后端；
  \item 给出正确性、健壮性、目标机器指令和性能证据；
  \item 形成可编译工程和完整 LaTeX 实验报告。
\end{enumerate}

\subsection{本实验完成内容}
\begin{table}[H]
\centering
\caption{工程实现范围}
\begin{tabularx}{\textwidth}{p{2.5cm}p{2.1cm}p{2.2cm}X}
\toprule
实现层次 & 后端 & 并行宽度 & 主要技术 \\
\midrule
正确性基线 & Scalar & 1 & 直接实现 $W/W'$、64 轮压缩和标准填充 \\
\midrule
x86-64 & AVX2 & 8 & YMM 8$\times$32-bit lane、移位拼接、向量 blend \\
x86-64 & AVX-512F & 16 & ZMM 16$\times$32-bit lane、\code{VPTERNLOGD}、\code{VPROLD}、opmask \\
ARM64 & NEON & 4 & \code{uint32x4\_t}、\code{EOR/ADD}、向量移位、\code{BSL} \\
\midrule
公共工程层 & Auto dispatch & 任意 & CPU 特性检测，宽向量优先，标量处理尾部 \\
\bottomrule
\end{tabularx}
\end{table}

\begin{keybox}{实验目标的准确边界}
本实现采用“多缓冲区并行”，即并行处理多条相互独立的 SM3 消息。它优化的是服务器批量验签、SM2-KDF、多连接握手、文件分片校验等场景的\textbf{吞吐率}，不能把 8 路或 16 路加速直接解释为单条消息延迟降低 8 倍或 16 倍。
\end{keybox}

\section{SM3 算法原理}

\subsection{总体结构}
SM3 接收长度小于 $2^{64}$ bit 的消息，按 512 bit 分组处理，输出 256 bit 杂凑值。设原始消息长度为 $l$，先附加一个比特 1，再附加 $k$ 个比特 0，使
\begin{equation}
  l+1+k \equiv 448 \pmod{512},
\end{equation}
最后附加 64 bit 的消息长度大端编码。每个 512 bit 分组被拆分为 16 个 32 bit 字，进行消息扩展和 64 轮压缩。

初始向量为：
\begin{align*}
IV={}&7380166f\;4914b2b9\;172442d7\;da8a0600\\
    &a96f30bc\;163138aa\;e38dee4d\;b0fb0e4e.
\end{align*}

\subsection{消息扩展}
对第 $i$ 个 512 bit 分组，首先解析得到 $W_0,\ldots,W_{15}$，随后计算
\begin{equation}
W_j=P_1\left(W_{j-16}\oplus W_{j-9}\oplus \rotl(W_{j-3},15)\right)
\oplus \rotl(W_{j-13},7)\oplus W_{j-6},\quad 16\le j<68,
\end{equation}
以及
\begin{equation}
W'_j=W_j\oplus W_{j+4},\quad 0\le j<64.
\end{equation}
其中
\begin{align}
P_0(X)&=X\oplus \rotl(X,9)\oplus \rotl(X,17),\\
P_1(X)&=X\oplus \rotl(X,15)\oplus \rotl(X,23).
\end{align}

\subsection{压缩函数}
设工作变量为 $A,B,C,D,E,F,G,H$。轮常量为
\begin{equation}
T_j=\begin{cases}
\texttt{0x79cc4519}, & 0\le j<16,\\
\texttt{0x7a879d8a}, & 16\le j<64.
\end{cases}
\end{equation}
布尔函数为
\begin{align}
FF_j(X,Y,Z)&=\begin{cases}X\oplus Y\oplus Z,&j<16,\\(X\land Y)\lor(X\land Z)\lor(Y\land Z),&j\ge16,\end{cases}\\
GG_j(X,Y,Z)&=\begin{cases}X\oplus Y\oplus Z,&j<16,\\(X\land Y)\lor(\neg X\land Z),&j\ge16.\end{cases}
\end{align}
每轮计算
\begin{align}
SS1&=\rotl\left(\rotl(A,12)+E+\rotl(T_j,j),7\right),\\
SS2&=SS1\oplus\rotl(A,12),\\
TT1&=FF_j(A,B,C)+D+SS2+W'_j,\\
TT2&=GG_j(E,F,G)+H+SS1+W_j,
\end{align}
然后更新工作变量，并在 64 轮结束后与输入链值异或。

\section{总体优化设计}

\subsection{为何采用多消息 SoA 并行}
单条 SM3 消息的 64 轮之间存在明显链式依赖：下一轮的 $A$ 和 $E$ 依赖本轮的 $TT1$ 和 $P_0(TT2)$，难以直接把一条消息的多轮放入不同 SIMD lane 并行执行。因此，本实验把不同消息映射到 lane，把同一轮的相同变量组成向量：
\begin{equation}
\vect{A}=(A_0,A_1,\ldots,A_{N-1}),\qquad
\vect{W_j}=(W_{j,0},W_{j,1},\ldots,W_{j,N-1}).
\end{equation}
这样，一条向量加法或异或指令可以同时完成 $N$ 条消息的同一操作。

\begin{figure}[H]
\centering
\begin{tikzpicture}[font=\small,>=Latex]
  \node[draw,rounded corners,fill=lightblue,minimum width=3.0cm,minimum height=1.0cm] (batch) {输入消息批次};
  \node[draw,rounded corners,fill=lightorange,minimum width=3.2cm,minimum height=1.0cm,right=1.0cm of batch,align=center] (gpr) {通用寄存器控制层\\指针/长度/索引/掩码};
  \node[draw,rounded corners,fill=lightgreen,minimum width=3.5cm,minimum height=1.0cm,right=1.0cm of gpr,align=center] (pack) {SoA 打包\\同一状态字进入同一向量};
  \node[draw,rounded corners,fill=purple!8,minimum width=3.8cm,minimum height=1.0cm,below=1.05cm of pack,align=center] (simd) {SIMD 压缩核心\\W/W'、P0/P1、FF/GG、模加};
  \node[draw,rounded corners,fill=lightorange,minimum width=3.2cm,minimum height=1.0cm,left=1.0cm of simd,align=center] (mask) {掩码提交\\只更新活跃 lane};
  \node[draw,rounded corners,fill=lightblue,minimum width=3.0cm,minimum height=1.0cm,left=1.0cm of mask] (out) {输出 256-bit 摘要};
  \draw[->,thick] (batch) -- (gpr);
  \draw[->,thick] (gpr) -- (pack);
  \draw[->,thick] (pack) -- (simd);
  \draw[->,thick] (simd) -- (mask);
  \draw[->,thick] (mask) -- (out);
  \draw[->,thick,dashed] (gpr.south) |- node[pos=0.3,right]{活跃位图} (mask.north);
\end{tikzpicture}
\caption{SIMD 与通用寄存器混合的总体数据流}
\end{figure}

\subsection{寄存器职责划分}
\begin{table}[H]
\centering
\caption{混合寄存器职责}
\begin{tabularx}{\textwidth}{p{3.2cm}X X}
\toprule
类别 & 保存内容 & 设计原因 \\
\midrule
SIMD 寄存器 & $A\sim H$、$W_j$、$W'_j$、轮中间量、活跃 lane 候选状态 & 这些变量在不同消息之间结构一致，可按 32-bit lane 并行 \\
通用寄存器 & 消息基址、长度、分组号、lane 循环、活跃位图、尾块数量 & 控制流和不规则地址计算不适合强行向量化 \\
栈/局部数组 & 68 个向量扩展字、每 lane 的 128 B 尾块、打包缓冲 & 降低 SIMD/GPR 之间频繁搬运，简化不同长度消息处理 \\
\bottomrule
\end{tabularx}
\end{table}

\subsection{不同长度消息的活跃 lane 机制}
对一批消息，最长消息可能包含更多完整分组。每处理一个分组序号 $b$，通用寄存器计算
\begin{equation}
active_i=[b < \lfloor len_i/64\rfloor].
\end{equation}
非活跃 lane 输入零分组，但压缩结果不会写回链值。三种后端的写回方式分别为：
\begin{itemize}
  \item AVX2：\code{\_mm256\_blendv\_epi8(old, candidate, mask)}；
  \item AVX-512：\code{\_mm512\_mask\_mov\_epi32(old, k, candidate)}；
  \item NEON：\code{vbslq\_u32(mask, candidate, old)}。
\end{itemize}
每条消息的尾部独立构造 1 或 2 个填充分组，因此可以正确覆盖 55/56、63/64、119/120 等边界。

\begin{figure}[H]
\centering
\begin{tikzpicture}[font=\small]
  \foreach \x/\lab in {0/msg0,1/msg1,2/msg2,3/msg3,4/msg4,5/msg5,6/msg6,7/msg7} {
    \draw[fill=lightblue] (\x*1.25,0) rectangle ++(1.05,0.72);
    \node at (\x*1.25+0.525,0.36) {\lab};
  }
  \node[anchor=east] at (-0.2,0.36) {$\vect{A}=$};
  \foreach \x in {0,1,2,3,4} {
    \draw[fill=lightgreen] (\x*1.25,-1.25) rectangle ++(1.05,0.72);
  }
  \foreach \x in {5,6,7} {
    \draw[fill=black!10] (\x*1.25,-1.25) rectangle ++(1.05,0.72);
  }
  \node[anchor=east] at (-0.2,-0.89) {active mask};
  \node[below] at (2.5,-1.35) {继续提交状态};
  \node[below] at (7.55,-1.35) {保持旧状态};
\end{tikzpicture}
\caption{AVX2 八路 SoA 状态与活跃 lane 掩码示意}
\end{figure}

\section{标量参考实现}

\subsection{实现作用}
标量版本不追求极致性能，而承担三个职责：
\begin{enumerate}
  \item 作为标准向量验证的最小可信实现；
  \item 作为 SIMD 随机差分测试的 oracle；
  \item 处理批量数量不足向量宽度时的尾部消息。
\end{enumerate}

\subsection{核心代码}
\begin{lstlisting}[style=cstyle,caption={标量消息扩展与压缩核心（节选）}]
for (unsigned j = 16; j < 68; ++j) {
    uint32_t x = w[j-16] ^ w[j-9] ^ ROTL32(w[j-3], 15);
    w[j] = P1(x) ^ ROTL32(w[j-13], 7) ^ w[j-6];
}
for (unsigned j = 0; j < 64; ++j) {
    uint32_t a12 = ROTL32(a, 12);
    uint32_t ss1 = ROTL32(a12 + e + T_ROT(j), 7);
    uint32_t ss2 = ss1 ^ a12;
    uint32_t tt1 = FF(j,a,b,c) + d + ss2 + (w[j] ^ w[j+4]);
    uint32_t tt2 = GG(j,e,f,g) + h + ss1 + w[j];
    d = c; c = ROTL32(b,9); b = a; a = tt1;
    h = g; g = ROTL32(f,19); f = e; e = P0(tt2);
}
\end{lstlisting}

\subsection{端序与填充}
SM3 以大端 32-bit 字解释输入。工程通过显式 \code{sm3\_load\_be32} 和 \code{sm3\_store\_be32} 避免主机端序依赖；尾块由统一函数构造，并写入 64-bit 大端比特长度。所有 SIMD 后端复用同一套端序和填充逻辑，降低平台实现之间出现语义分叉的风险。

\section{x86 AVX2 八路实现}

\subsection{数据布局}
AVX2 的 YMM 寄存器宽度为 256 bit，可容纳 8 个 32-bit lane。本实现将 8 条消息的同一状态字放入一个 \code{\_\_m256i}：
\begin{align*}
\vect{A}&=(A_0,A_1,\ldots,A_7),\\
\vect{W_j}&=(W_{j,0},W_{j,1},\ldots,W_{j,7}).
\end{align*}
因此，\code{VPADDD}、\code{VPXOR}、\code{VPSLLD/VPSRLD} 可同步完成八条消息的轮操作。

\subsection{循环移位与置换函数}
AVX2 基础整数指令没有统一的按 32-bit 立即数循环左移 intrinsic，本实现用左右移和或组合：
\begin{lstlisting}[style=cstyle,caption={AVX2 循环左移和 $P_1$}]
static inline __m256i rotl32(__m256i x, int n) {
    return _mm256_or_si256(_mm256_slli_epi32(x, n),
                           _mm256_srli_epi32(x, 32 - n));
}
static inline __m256i p1(__m256i x) {
    return _mm256_xor_si256(x,
           _mm256_xor_si256(rotl32(x,15), rotl32(x,23)));
}
\end{lstlisting}

\subsection{布尔函数}
前 16 轮的 $FF/GG$ 均为三输入异或；后 48 轮分别为 majority 和 choose。本实现使用 \code{AND/OR/ANDNOT/XOR} 组合，保持所有中间量位于 YMM 寄存器中。

\subsection{不同长度 lane 写回}
压缩函数先计算候选链值
\begin{equation}
\vect{V}^{candidate}=\vect{V}^{old}\oplus(A,B,C,D,E,F,G,H),
\end{equation}
再通过全 0/全 1 的 lane 掩码执行 \code{blendv}。这样，即使某些消息已完成，其他消息仍可继续在同一个批次中处理后续分组。

\subsection{AVX-SSE 状态清理}
函数结束调用 \code{\_mm256\_zeroupper()}，避免后续执行传统 SSE 代码时产生不必要的 AVX-SSE 状态转换代价。

\section{x86 AVX-512 十六路实现}

\subsection{向量宽度与掩码寄存器}
AVX-512F 的 ZMM 寄存器宽度为 512 bit，可容纳 16 个 32-bit lane。与 AVX2 相比，本实现不仅把批量宽度翻倍，还使用 \code{\_\_mmask16} 保存活跃 lane，在状态写回阶段直接执行 masked move。

\subsection{三元逻辑优化}
\code{VPTERNLOGD} 可用 8-bit 真值表编码任意三输入布尔函数。本实现采用：
\begin{table}[H]
\centering
\caption{AVX-512 三元逻辑立即数}
\begin{tabular}{lll}
\toprule
函数 & 逻辑表达式 & \code{imm8} \\
\midrule
XOR3 & $X\oplus Y\oplus Z$ & \code{0x96} \\
Majority & $(X\land Y)\lor(X\land Z)\lor(Y\land Z)$ & \code{0xE8} \\
Choose & $(X\land Y)\lor(\neg X\land Z)$ & \code{0xCA} \\
\bottomrule
\end{tabular}
\end{table}
这使 $P_0/P_1$ 的最终三输入异或、前 16 轮 $FF/GG$ 以及后 48 轮布尔函数均可减少逻辑指令数量。

\begin{lstlisting}[style=cstyle,caption={AVX-512 三元逻辑实现}]
static inline __m512i xor3(__m512i x, __m512i y, __m512i z) {
    return _mm512_ternarylogic_epi32(x, y, z, 0x96);
}
static inline __m512i majority(__m512i x, __m512i y, __m512i z) {
    return _mm512_ternarylogic_epi32(x, y, z, 0xe8);
}
static inline __m512i choose(__m512i x, __m512i y, __m512i z) {
    return _mm512_ternarylogic_epi32(x, y, z, 0xca);
}
\end{lstlisting}

\subsection{循环移位}
编译器在 AVX-512 路径中能够把固定计数的 32-bit 循环移位识别为 \code{VPROLD}。反汇编中可观察到例如 \code{vprold zmm4,zmm4,0xf}，证明目标代码没有停留在伪向量化源代码层面。

\subsection{频率与平台权衡}
AVX-512 的理论 lane 数是 AVX2 的两倍，但实际性能不一定严格翻倍，原因包括执行端口、寄存器压力、栈上扩展字读写、内存带宽以及部分处理器的宽向量降频。报告因此同时给出吞吐率、加速比和实验平台，不用向量宽度直接替代实测结果。

\section{ARM64 NEON 四路实现}

\subsection{架构特点}
ARM64 提供 32 个 64-bit 通用寄存器和 32 个 128-bit SIMD/FP 寄存器。NEON 的 \code{uint32x4\_t} 可同时处理四个 32-bit lane。本实验沿用与 x86 相同的 SoA 语义，以保证测试模型和代码结构一致。

\subsection{循环移位}
NEON intrinsic 版本使用带符号移位计数的 \code{vshlq\_u32}：正计数左移，负计数右移，再使用 \code{vorrq\_u32} 合并。
\begin{lstlisting}[style=cstyle,caption={ARM64 NEON 循环左移}]
static inline uint32x4_t rotl32(uint32x4_t x, int n) {
    int32x4_t left  = vdupq_n_s32(n);
    int32x4_t right = vdupq_n_s32(n - 32);
    return vorrq_u32(vshlq_u32(x, left),
                     vshlq_u32(x, right));
}
\end{lstlisting}

\subsection{掩码提交}
NEON 通过 \code{vbslq\_u32(mask, candidate, old)} 按位选择。活跃 lane 的掩码为 \code{0xffffffff}，非活跃 lane 为 0，因此可以在不引入每 lane 分支的情况下保持已完成消息的链值。

\subsection{交叉编译验证与运行边界}
当前报告生成环境为 x86-64，无法直接执行 AArch64 二进制。本工程已将 \code{sm3\_neon.c} 交叉编译为 AArch64 ELF 可重定位对象，并在反汇编中确认出现 \code{dup v*.4s}、\code{bsl} 等 NEON 指令。这证明源代码能够生成目标 ISA；但 ARM64 正确性和性能仍应在真实设备上运行同一测试程序。本报告不把课程材料或模拟数据冒充本工程实测数据。

\begin{table}[H]
\centering
\caption{课程材料中的 ARM64 参考性能（非本工程实测）}
\begin{tabular}{lrrrrl}
\toprule
架构 & 标准 MB/s & 标准 CPB & 优化 MB/s & 优化 CPB & 推荐实现 \\
\midrule
Cortex-A55 & 65.04 & 27.78 & 121.12 & 14.92 & 通用寄存器 \\
Cortex-A76 & 121.89 & 16.46 & 293.94 & 7.99 & 混合寄存器 \\
ARM A78/Android & 168.50 & 12.79 & 312.97 & 6.89 & 混合寄存器 \\
\bottomrule
\end{tabular}
\end{table}
课程数据说明：不同 ARM 微架构上的最优策略不完全相同。部分低功耗核心的 NEON 吞吐能力有限，全通用寄存器实现可能更合适；较强核心通常能从混合并行中获益。

\section{运行时分发与工程结构}

\subsection{后端选择}
公共函数 \code{sm3\_hash\_many} 接收后端枚举、消息指针数组、长度数组和输出数组。自动模式的优先级为：
\begin{equation}
\text{x86: AVX-512(16)}\rightarrow\text{AVX2(8)}\rightarrow\text{Scalar(1)},
\end{equation}
\begin{equation}
\text{ARM64: NEON(4)}\rightarrow\text{Scalar(1)}.
\end{equation}
例如 37 条 x86 消息在 AVX-512 可用时被拆分为 $16+16+5$，最后 5 条由标量处理；在只有 AVX2 时被拆分为 $8+8+8+8+5$。

\begin{lstlisting}[style=cstyle,caption={x86 自动分发逻辑（简化）}]
if (cpu_has_avx512f()) {
    while (remaining >= 16) hash16_avx512(...);
    if (remaining >= 8 && cpu_has_avx2()) hash8_avx2(...);
    scalar_tail(...);
} else if (cpu_has_avx2()) {
    while (remaining >= 8) hash8_avx2(...);
    scalar_tail(...);
} else {
    scalar_all(...);
}
\end{lstlisting}

\subsection{避免非法指令}
\code{sm3\_avx2.c} 和 \code{sm3\_avx512.c} 分别使用 \code{-mavx2} 和 \code{-mavx512f} 单独编译，调度和标量文件不启用高级 ISA。这样，程序可先在普通指令路径启动并检查 CPU/OS 支持，再调用对应函数；若对整个工程统一使用 \code{-mavx512f}，编译器可能在初始化或普通代码中生成 AVX-512，导致不支持的 CPU 直接触发非法指令。

\subsection{工程目录}
{\footnotesize
\begin{longtable}{p{5.2cm}p{8.8cm}}
\caption{项目文件说明}\\
\toprule
文件 & 作用 \\
\midrule
\endfirsthead
\toprule
文件 & 作用 \\
\midrule
\endhead
\nolinkurl{include/sm3.h} & 公共 API、后端枚举、lane 宽度和 CPU 检测声明 \\
\nolinkurl{src/sm3_internal.h} & 初始向量、端序、循环移位、$P_0/P_1$、轮常量和尾块构造 \\
\nolinkurl{src/sm3_scalar.c} & 单消息标量参考实现 \\
\nolinkurl{src/sm3_avx2.c} & AVX2 八路多消息实现 \\
\nolinkurl{src/sm3_avx512.c} & AVX-512 十六路多消息实现 \\
\nolinkurl{src/sm3_neon.c} & ARM64 NEON 四路多消息实现 \\
\nolinkurl{src/sm3_dispatch.c} & CPU 特性检测与任意数量消息的自动分发 \\
\nolinkurl{tests/test_sm3.c} & 标准向量、边界、随机差分和调度测试 \\
\nolinkurl{bench/bench_sm3.c} & 预热、7 次采样、中位数吞吐率基准 \\
\nolinkurl{tools/sm3_cli.c} & 文件摘要工具，供 OpenSSL 差分脚本调用 \\
\nolinkurl{scripts/verify_openssl.py} & 生成边界/随机文件并逐字节比较 OpenSSL 摘要 \\
\nolinkurl{scripts/run_validation.sh} & 一键收集环境、测试、性能、Sanitizer 和反汇编证据 \\
\bottomrule
\end{longtable}
}

\section{正确性与健壮性验证}

\subsection{标准向量}
测试程序验证了以下公开标准用例：
\begin{table}[H]
\centering
\caption{SM3 标准向量}
\begin{tabularx}{\textwidth}{p{3.1cm}X}
\toprule
消息 & 期望摘要 \\
\midrule
空消息 & \code{1ab21d8355cfa17f8e61194831e81a8f22bec8c728fefb747ed035eb5082aa2b} \\
\code{abc} & \code{66c7f0f462eeedd9d1f2d46bdc10e4e24167c4875cf2f7a2297da02b8f4ba8e0} \\
\code{abcd} 重复 16 次 & \code{debe9ff92275b8a138604889c18e5a4d6fdb70e5387e5765293dcba39c0c5732} \\
\bottomrule
\end{tabularx}
\end{table}

\subsection{边界与随机差分}
边界长度覆盖：
\begin{equation*}
0,1,2,3,7,15,31,55,56,57,63,64,65,119,120,121,127,128,129,255,256,1024,8192.
\end{equation*}
测试对每个 SIMD lane 生成不同长度和不同内容的消息，并与标量 \code{sm3\_digest} 逐字节比较。当前验证输出为：
\begin{verbatim}
[PASS] avx2-8way     edge/random differential (101 rounds, 8 messages)
[PASS] avx512-16way  edge/random differential (101 rounds, 16 messages)
[PASS] auto          edge/random differential (51 rounds, 37 messages)
All SM3 tests passed.
\end{verbatim}

\subsection{OpenSSL 差分}
脚本生成固定边界文件和随机文件，通过工程 CLI 与 \code{openssl dgst -sm3 -binary} 分别计算摘要，结果为：
\begin{verbatim}
[PASS] 128 files matched OpenSSL SM3; backend auto-dispatch exercised
\end{verbatim}
这项测试同时验证了文件输入、任意长度处理、运行时分发和摘要端序。

\subsection{Sanitizer}
使用以下选项重新构建：
\begin{lstlisting}[style=shellstyle,caption={ASan/UBSan 验证命令}]
make sanitize
# -fsanitize=address,undefined
# -fno-omit-frame-pointer
\end{lstlisting}
标准向量和全部随机差分仍通过，输出中未出现 heap/stack buffer overflow、use-after-free 或 undefined behavior 报告。

\section{目标指令证据}

\subsection{x86 AVX2}
反汇编中观察到：
\begin{verbatim}
vmovdqa ymm0, YMMWORD PTR [...]
vpsrld  ymm2, ymm4, 0x11
vpxor   ymm1, ymm6, ymm10
vpslld  ymm4, ymm4, 0xf
\end{verbatim}
这证明核心路径使用了 YMM 寄存器和按 32-bit lane 的向量逻辑/移位。

\subsection{x86 AVX-512}
反汇编中观察到：
\begin{verbatim}
vmovdqa64 zmm1, ZMMWORD PTR [...]
vprold    zmm4, zmm4, 0xf
vpternlogd ...
\end{verbatim}
测试主机真实支持 AVX-512F，故 AVX-512 路径不仅完成编译，也被运行时测试和 Benchmark 实际执行。

\subsection{ARM64 NEON}
交叉编译对象信息为：
\begin{verbatim}
ELF 64-bit LSB relocatable, ARM aarch64, version 1 (SYSV)
\end{verbatim}
反汇编中观察到：
\begin{verbatim}
dup v5.4s, w8
bsl v2.16b, v1.16b, v0.16b
\end{verbatim}
说明编译器生成了 128-bit NEON lane 操作和掩码选择指令。

\section{性能测试设计}

\subsection{测试原则}
基准程序采用以下方法降低噪声：
\begin{enumerate}
  \item 计时前执行 3 次预热；
  \item 每个后端执行相同的 16 条消息工作量；
  \item 随机数据生成、内存分配和打印不计入计时区间；
  \item 每个后端独立采样 7 次并取中位数；
  \item 使用 volatile sink 消费摘要，阻止编译器删除计算；
  \item 同时报告 MiB/s、ns/B 和相对标量加速比。
\end{enumerate}

\subsection{公平比较方式}
对 16 条消息：
\begin{itemize}
  \item Scalar：顺序计算 16 次；
  \item AVX2：执行两批八路；
  \item AVX-512：执行一批十六路；
  \item ARM64 NEON：执行四批四路。
\end{itemize}
因此，各后端处理的总字节数一致。

\section{x86 性能结果与分析}

\subsection{本人 AMD Ryzen 5 5600U / WSL2 实测}
以下数据来自本人已保存的固定 CPU 核测试输出。该主机支持 AVX2，不支持 AVX-512，因此 AVX-512 被正确跳过。

\begin{table}[H]
\centering
\caption{本人 x86 实测：标量与 AVX2 批量吞吐率}
\begin{tabular}{rrrrr}
\toprule
每条消息长度 & Scalar (MiB/s) & AVX2 (MiB/s) & AVX2 ns/B & 加速比 \\
\midrule
64 B & 62.42 & 534.53 & 1.784 & 8.56$\times$ \\
256 B & 145.70 & 972.65 & 0.980 & 6.68$\times$ \\
1024 B & 176.17 & 1135.64 & 0.840 & 6.45$\times$ \\
4096 B & 183.32 & 1256.63 & 0.759 & 6.85$\times$ \\
16384 B & 186.95 & 1174.12 & 0.812 & 6.28$\times$ \\
1 MiB & 195.23 & 1320.00 & 0.722 & 6.76$\times$ \\
\bottomrule
\end{tabular}
\end{table}

\begin{figure}[H]
\centering
\begin{tikzpicture}
\begin{axis}[
  width=0.91\textwidth,height=7.0cm,
  xlabel={每条消息长度（字节，对数坐标）},
  ylabel={相对标量加速比},
  xmode=log, log basis x=2,
  ymin=0,ymax=10,
  grid=both,
  major grid style={black!15},
  minor grid style={black!7},
  xtick={64,256,1024,4096,16384,1048576},
  xticklabels={64,256,1024,4096,16384,1 MiB},
  tick label style={font=\small},
  label style={font=\small},
]
\addplot[very thick,mark=*,deepblue] coordinates {
  (64,8.56) (256,6.68) (1024,6.45) (4096,6.85) (16384,6.28) (1048576,6.76)
};
\end{axis}
\end{tikzpicture}
\caption{本人 x86 环境下 AVX2 相对标量加速比}
\end{figure}

\subsection{独立 AVX-512 验证环境}
为了实际执行 AVX-512 路径，工程还在一台支持 AVX-512F 的 x86-64 验证环境中运行。环境为 AMD EPYC 9V74 虚拟 CPU、GCC 14.2.0、OpenSSL 3.5.5。4096 B、16 条消息、7 次采样中位数如下。

\begin{table}[H]
\centering
\caption{AVX-512 验证环境的 4096 B 批量性能}
\begin{tabular}{lrrr}
\toprule
后端 & MiB/s & ns/B & 相对标量 \\
\midrule
Scalar & 143.042 & 6.6671 & 1.000$\times$ \\
AVX2 8-way & 1034.520 & 0.9219 & 7.232$\times$ \\
AVX-512 16-way & 1843.984 & 0.5172 & 12.891$\times$ \\
\bottomrule
\end{tabular}
\end{table}

\begin{figure}[H]
\centering
\begin{tikzpicture}
\begin{axis}[
  ybar,
  width=0.82\textwidth,height=7cm,
  ylabel={吞吐率（MiB/s）},
  symbolic x coords={Scalar,AVX2,AVX-512},
  xtick=data,
  ymin=0,ymax=2100,
  nodes near coords,
  nodes near coords align={vertical},
  bar width=28pt,
  grid=major,
  major grid style={black!12},
]
\addplot[fill=deepblue!55,draw=deepblue] coordinates {(Scalar,143.042) (AVX2,1034.520) (AVX-512,1843.984)};
\end{axis}
\end{tikzpicture}
\caption{独立验证环境中不同 x86 后端吞吐率}
\end{figure}

\subsection{结果解释}
\begin{enumerate}
  \item AVX2 的实际加速约为 6--8 倍，与 8 路 lane 宽度接近，但受打包、尾块、掩码、内存访问和指令依赖影响，不会始终等于 8 倍。
  \item 64 B 消息中一次调用包含大量独立短消息，标量函数调用和控制开销占比较高，因此批量 SIMD 的相对优势可能更突出。
  \item AVX-512 在验证环境中达到约 12.89 倍，而不是严格 16 倍，说明 ZMM 宽度只是上限，实际还受扩展字存储、端口吞吐和平台频率策略影响。
  \item 不同机器上的绝对 MiB/s 不能直接横向比较；应优先比较同一机器、同一编译器、同一工作量下的后端差异。
\end{enumerate}

\section{ARM64 测试与复测方案}

\subsection{真实设备命令}
在 ARM64 Linux 上直接运行：
\begin{lstlisting}[style=shellstyle,caption={ARM64 原生测试命令}]
make clean
make -j"$(nproc)" all
make test
make verify-openssl
make asm-check
./build/sm3_bench 64 256
./build/sm3_bench 256 256
./build/sm3_bench 1024 256
./build/sm3_bench 4096 256
./build/sm3_bench 16384 256
\end{lstlisting}

\subsection{建议采集指标}
除 MiB/s 和 ns/B 外，建议在原生 Linux 使用：
\begin{lstlisting}[style=shellstyle,caption={ARM64 perf 采集示例}]
taskset -c 2 perf stat -r 10 \
  -e cycles,instructions,branches,branch-misses,cache-references,cache-misses \
  ./build/sm3_bench 4096 512
\end{lstlisting}
计算
\begin{equation}
\mathrm{cycles/byte}=\frac{cycles}{\text{实际处理总字节数}},\qquad
\mathrm{instructions/byte}=\frac{instructions}{\text{实际处理总字节数}}.
\end{equation}

\section{问题、限制与改进方向}

\subsection{当前限制}
\begin{enumerate}
  \item 当前 ARM64 路径完成了交叉编译和反汇编验证，但报告生成主机不是 ARM64，尚缺真实设备上的运行数据；
  \item SIMD 路径当前先构造完整的 $W_0\sim W_{67}$ 向量数组，栈读写较多；
  \item 输入从 AoS 指针数组打包为 SoA 向量时使用标量 gather，可能成为短消息开销；
  \item 实现侧重批量吞吐率，未针对单消息在线消息扩展或四轮展开做极致延迟优化；
  \item Benchmark 未统一控制所有机器的睿频、温度和虚拟化调度。
\end{enumerate}

\subsection{进一步优化}
\begin{enumerate}
  \item \textbf{消息扩展与压缩交错：}按需生成 $W_j$，减少 68 个向量的栈占用和访存；
  \item \textbf{四轮展开：}利用状态字四轮后回到对应寄存器位置，减少显式数据移动；
  \item \textbf{AVX-512 掩码加载：}用 masked load/gather 降低 lane 打包分支；
  \item \textbf{ARM64 微架构适配：}为 Cortex-A55 提供偏通用寄存器版本，为 A76/A78 提供混合版本；
  \item \textbf{固定长度专用入口：}对 SM2-KDF 等固定布局场景预计算公共分组并批量处理计数器；
  \item \textbf{更严格的基准：}固定核心、关闭不必要服务、记录温度和频率、使用多次独立进程中位数；
  \item \textbf{持续集成：}在 x86 AVX2、x86 AVX-512 和 ARM64 三类 runner 上自动执行标准向量、OpenSSL 差分和 Sanitizer。
\end{enumerate}

\section{实验结论}
本实验完成了 SM3 标量实现，以及 x86 AVX2 八路、x86 AVX-512 十六路和 ARM64 NEON 四路多消息优化实现。工程没有简单地“把所有变量塞进 SIMD”，而是按规则划分职责：规则、密集、同构的压缩运算使用 SIMD 寄存器；消息指针、长度、分组调度、活跃 lane 和尾块构造使用通用寄存器。不同长度消息通过掩码提交机制保持正确状态，任意消息数量通过宽向量到窄向量再到标量的分层调度完成。

验证结果形成了从算法到机器码的证据闭环：标准向量、边界和随机差分全部通过；128 个文件与 OpenSSL SM3 一致；ASan/UBSan 未报告错误；反汇编确认 YMM、ZMM、\code{VPROLD/VPTERNLOGD} 以及 ARM64 NEON \code{BSL} 等目标指令。本人 x86 实测中，AVX2 在多种长度上取得约 6.28--8.56 倍加速；支持 AVX-512 的验证环境中，AVX-512 在 4096 B 消息上取得约 12.89 倍加速。

综合来看，本实验满足“ARM64 与 x86、AVX2 与 AVX-512、SIMD 与通用寄存器混合”的作业要求，同时明确区分了源码可编译、目标指令存在、运行正确和性能实测四类证据。后续只需在真实 ARM64 设备上运行现成测试与基准命令，即可补齐该平台的实测性能表。

\clearpage
\appendix
\section{编译、测试与跑分命令}

\subsection{x86-64}
\begin{lstlisting}[style=shellstyle]
make clean
make -j"$(nproc)" all
make test
make verify-openssl
make asm-check
./build/sm3_bench 4096 256
make sanitize
\end{lstlisting}

\subsection{ARM64 原生}
\begin{lstlisting}[style=shellstyle]
make clean
make -j"$(nproc)" all
make test
make asm-check
./build/sm3_bench 4096 256
\end{lstlisting}

\subsection{x86 主机交叉编译 ARM64}
\begin{lstlisting}[style=shellstyle]
sudo apt install gcc-aarch64-linux-gnu binutils-aarch64-linux-gnu
make arm64
\end{lstlisting}

\subsection{一键验收}
\begin{lstlisting}[style=shellstyle]
TOTAL_MIB=256 ./scripts/run_validation.sh
\end{lstlisting}

\section{关键公共 API}
\begin{lstlisting}[style=cstyle]
typedef enum sm3_backend {
    SM3_BACKEND_SCALAR,
    SM3_BACKEND_AVX2,
    SM3_BACKEND_AVX512,
    SM3_BACKEND_NEON,
    SM3_BACKEND_AUTO
} sm3_backend;

void sm3_digest(const uint8_t *msg, size_t len, uint8_t out[32]);

int sm3_hash_many(sm3_backend backend,
                  const uint8_t *const *msg,
                  const size_t *len,
                  uint8_t (*out)[32],
                  size_t count);
\end{lstlisting}

\begin{thebibliography}{99}
\bibitem{GBT32905}
国家质量监督检验检疫总局、国家标准化管理委员会. GB/T 32905-2016 信息安全技术 SM3 密码杂凑算法, 2016.

\bibitem{SM3IETF}
IETF. The SM3 Cryptographic Hash Function. Internet-Draft, algorithm description corresponding to GB/T 32905-2016.

\bibitem{IntelGuide}
Intel Corporation. Intel Intrinsics Guide and Intel 64 and IA-32 Architectures Software Developer's Manual.

\bibitem{GCCx86}
GNU Project. GCC x86 Options: AVX2 and AVX-512 feature flags.

\bibitem{ArmNeon}
Arm Limited. Arm Neon Intrinsics Reference and Cortex-A Series Programmer's Guide for Armv8-A.

\bibitem{CourseSM3}
山东大学网络空间安全学院. SM3 软件实现课程讲义, 2026.
\end{thebibliography}

\end{document}
